\documentclass[a4paper]{ctexart}

\usepackage{amsmath}
\usepackage{amssymb}
\usepackage{amsthm}
\usepackage[top=2.6cm,bottom=2.8cm,left=2.6cm,right=2.6cm]{geometry}
\usepackage[colorlinks=true,linkcolor=blue,pdfpagelabels=false]{hyperref}
\usepackage{bookmark}

\linespread{1.25}

\title{Algebra III 全书考前复习}
\author{}
\date{2026-06-29}

\begin{document}
\maketitle
\tableofcontents

\section*{第0章：总览}
\phantomsection
\addcontentsline{toc}{section}{第0章：总览}

这门课的主线是：从谱与局部化出发，建立 Noether 环、整扩张、分次与完备化、维数理论，再进入
Cohen--Macaulay、normal、regular 这些最核心的交换代数对象。

如果只记一条总路线，可以记成：
\[
\text{Spec 与局部化}
\to
\text{flat / integral}
\to
\text{Noether / Artin / primary decomposition},
\]
\[
\text{completion 与 Hilbert 函数}
\to
\text{dimension}
\to
\text{CM / normal / regular}.
\]

\section{第1章：Rings and Modules}

\subsection{1.1 Prime ideals}

\paragraph{核心对象}
\begin{itemize}
  \item 素理想 \(\mathfrak p\): \(A/\mathfrak p\) 是整环。
  \item 极大理想 \(\mathfrak m\): \(A/\mathfrak m\) 是域。
  \item 谱空间 \(\operatorname{Spec}A\): 所有素理想组成的集合。
\end{itemize}

\paragraph{必须掌握的事实}
\begin{itemize}
  \item 任意非零环至少有一个极大理想。
  \item 任意环同态 \(f:A\to B\) 诱导连续映射
  \[
    f^*:\operatorname{Spec}B\to \operatorname{Spec}A,\qquad
    \mathfrak q\mapsto f^{-1}(\mathfrak q).
  \]
  \item 幂零根
  \[
    \operatorname{nilrad}(A)=\bigcap_{\mathfrak p\in \operatorname{Spec}A}\mathfrak p.
  \]
  \item 理想根
  \[
    \sqrt I=\bigcap_{\mathfrak p\supset I}\mathfrak p.
  \]
\end{itemize}

\paragraph{Zariski 拓扑}
闭集定义为
\[
V(I)=\{\mathfrak p\in \operatorname{Spec}A\mid I\subset \mathfrak p\}.
\]
要会立刻写出：
\[
V(0)=\operatorname{Spec}A,\qquad V(1)=\varnothing,\qquad
V(IJ)=V(I)\cup V(J).
\]

\paragraph{考前提示}
\begin{itemize}
  \item 交换代数里很多命题本质上都在讨论 \(\operatorname{Spec}A\) 上的几何行为。
  \item “闭点”通常对应极大理想。
  \item \(\operatorname{Spec}(A/I)\) 应直接想到 \(V(I)\)。
\end{itemize}

\subsection{1.2 Nakayama lemma}

\paragraph{Jacobson 根与局部环}
\begin{itemize}
  \item Jacobson 根
  \[
    J(A)=\bigcap_{\mathfrak m\in \operatorname{MaxSpec}A}\mathfrak m.
  \]
  \item 局部环就是只有唯一极大理想的环。
\end{itemize}

\paragraph{Nakayama 引理}
这是全书最常用工具之一。常见形式是：若 \(M\) 为有限生成 \(A\)-模，\(\mathfrak a\subset J(A)\)，且
\[
\mathfrak aM=M,
\]
则 \(M=0\)。

\paragraph{最常用推论}
\begin{itemize}
  \item 在局部环 \((A,\mathfrak m)\) 中，若 \(x_1,\dots,x_r\) 在 \(M/\mathfrak mM\) 中生成，则它们生成 \(M\)。
  \item \(M/\mathfrak mM\) 的 \(k=A/\mathfrak m\)-维数就是 \(M\) 的最少生成元个数。
\end{itemize}

\paragraph{考前提示}
很多“最少生成元个数”的结论，本质上都靠 Nakayama。

\section{第2章：Localization}

\subsection{2.1--2.2 局部化}

\paragraph{要点}
\begin{itemize}
  \item 给定乘法集 \(S\subset A\)，构造 \(S^{-1}A\)、\(S^{-1}M\)。
  \item 最重要的特例是
  \[
    A_{\mathfrak p}=(A\setminus \mathfrak p)^{-1}A.
  \]
  \item \(A_{\mathfrak p}\) 是局部环，极大理想是 \(\mathfrak pA_{\mathfrak p}\)。
\end{itemize}

\paragraph{谱的行为}
\[
\operatorname{Spec}(S^{-1}A)\cong \{\mathfrak p\in \operatorname{Spec}A\mid \mathfrak p\cap S=\varnothing\}.
\]
特别地，
\[
\operatorname{Spec}(A_{\mathfrak p})\cong \{\mathfrak q\in \operatorname{Spec}A\mid \mathfrak q\subset \mathfrak p\}.
\]

\paragraph{模论要点}
\begin{itemize}
  \item 局部化是正合函子。
  \item \(S^{-1}A\) 是 flat 的 \(A\)-代数。
  \item 模为零可以局部检验：
  \[
    M=0 \Longleftrightarrow M_{\mathfrak m}=0 \text{ 对所有极大理想 } \mathfrak m.
  \]
\end{itemize}

\paragraph{考前提示}
交换代数里的“局部化后看”几乎是默认动作。很多全局命题先化成局部命题，再在局部环里证明。

\section{第3章：Flatness}

\subsection{3.1 Flat modules}

\paragraph{定义}
\(M\) 是 flat，当且仅当函子
\[
-\otimes_A M
\]
是正合函子。等价地说：任意 \(A\)-模短正合列在张量 \(M\) 后仍保持正合。

\paragraph{等价刻画}
\begin{itemize}
  \item flat 等价于 \(\operatorname{Tor}_1^A(M,N)=0\) 对所有 \(N\)。
  \item 也可用有限生成理想的张量单射性来判别。
  \item flat 性可以局部检验：
  \[
    M \text{ flat over } A
    \Longleftrightarrow
    M_{\mathfrak p} \text{ flat over } A_{\mathfrak p}\ \forall \mathfrak p.
  \]
\end{itemize}

\paragraph{有限生成情形}
有限生成 flat 模在局部环上是自由模；这让 flat 和“局部自由”联系起来。

\subsection{3.2 Faithful flatness}

\paragraph{定义}
faithfully flat 比 flat 更强：flat 只要求张量函子保持正合；faithfully flat 则还要求它反映正合性。常用判据是：若 \(N\neq 0\)，则 \(N\otimes_A M\neq 0\)。

\paragraph{常用判别}
\begin{itemize}
  \item \(M\) faithfully flat
  \(\Longleftrightarrow\)
  \(M\) flat 且对任何非零模 \(N\)，有 \(N\otimes_A M\neq 0\)。
  \item 对环同态 \(A\to B\)，
  \[
    B \text{ faithfully flat over } A
    \Longleftrightarrow
    B \text{ flat 且 } \operatorname{Spec}B\to \operatorname{Spec}A \text{ 满射}.
  \]
\end{itemize}

\subsection{3.3 Flat ring homomorphisms}

\paragraph{最重要结论}
flat 环同态满足 \textbf{going-down}。

\paragraph{考前提示}
\begin{itemize}
  \item integral 扩张给 going-up。
  \item flat 扩张给 going-down。
  \item 这两个方向在后续维数理论中反复出现。
\end{itemize}

\section{第4章：Integral extensions}

\subsection{4.1 Integral dependence, going-up, going-down}

\paragraph{定义}
元素 \(x\in B\) 对 \(A\) 是整的，指满足首一方程
\[
x^n+a_1x^{n-1}+\cdots+a_n=0,\qquad a_i\in A.
\]

\paragraph{必须会的结论}
\begin{itemize}
  \item 若 \(B\) 是有限生成 \(A\)-模，则 \(B\) 的每个元素都整于 \(A\)。
  \item integral closure 在扩张中形成子环。
  \item 整性对商与局部化稳定。
\end{itemize}

\paragraph{谱上的性质}
若 \(A\subset B\) 为 integral 扩张，则
\begin{itemize}
  \item \(\operatorname{Spec}B\to \operatorname{Spec}A\) 满射；
  \item lying-over 成立；
  \item going-up 成立。
\end{itemize}

\paragraph{整闭}
整环 \(A\) 在分式域中整闭，称为 integrally closed 或 normal domain。
要会记住：
\[
A \text{ 整闭}
\Longleftrightarrow
A_{\mathfrak p} \text{ 对所有 } \mathfrak p \text{ 整闭}.
\]

\paragraph{going-down}
若 \(A\subset B\) 是 integral 扩张且 \(A\) 整闭，则 going-down 成立。

\subsection{4.2 Integral closures and valuation rings}

\paragraph{valuation ring}
对分式域 \(K\) 中的子环 \(B\)，若任意 \(0\neq x\in K\) 满足 \(x\in B\) 或 \(x^{-1}\in B\)，则 \(B\) 是 valuation ring。

\paragraph{关键思想}
整闭包可以写成所有包含原环的 valuation rings 的交。

\paragraph{意义}
valuation ring 是“尽可能线性排序”的局部模型，常用来检测整闭性。

\subsection{4.3 Noether normalization and Nullstellensatz}

\paragraph{Noether normalization}
若 \(A\) 是有限生成 \(k\)-代数，则存在代数无关元 \(t_1,\dots,t_d\)，使得
\[
A \text{ integral over } k[t_1,\dots,t_d].
\]

\paragraph{弱零点定理}
若 \(A\) 是有限生成 \(k\)-代数且是域，则在代数闭域情形下 \(A\) 其实就是 \(k\)。

\paragraph{Hilbert Nullstellensatz}
最大理想与几何点对应，是代数几何和交换代数接轨的核心定理。

\paragraph{考前提示}
\begin{itemize}
  \item Noether normalization 是“有限生成代数有限覆盖一个多项式环”。
  \item 它直接导向后面的超越维数公式。
\end{itemize}

\section{第5章：Noetherian and Artinian rings}

\subsection{5.1 Chain conditions}

\paragraph{定义}
\begin{itemize}
  \item Noetherian: 满足升链条件。
  \item Artinian: 满足降链条件。
  \item finite length: 有 composition series。
\end{itemize}

\paragraph{必须会的结论}
\begin{itemize}
  \item 模 Noetherian 当且仅当每个子模有限生成。
  \item finite length 等价于同时 Noetherian 与 Artinian。
  \item composition series 的长度唯一，这就是 Jordan--H\"older 型结论。
\end{itemize}

\subsection{5.2 Noetherian rings}

\paragraph{最关键判别}
\[
A \text{ Noetherian}
\Longleftrightarrow
\text{每个理想有限生成}.
\]

\paragraph{稳定性}
\begin{itemize}
  \item 有限生成模、商环、有限模扩张、局部化都保持 Noetherian。
  \item Hilbert basis theorem:
  \[
    A \text{ Noetherian} \Rightarrow A[X] \text{ Noetherian}.
  \]
  因而
  \[
    A[X_1,\dots,X_n] \text{ Noetherian}.
  \]
\end{itemize}

\subsection{5.3 Artinian rings}

\paragraph{核心结论}
\begin{itemize}
  \item Artinian 环中每个素理想都是极大理想。
  \item
  \[
    J(A)=\operatorname{nilrad}(A),
  \]
  且该理想幂零。
  \item Artinian 环只有有限多个极大理想。
  \item
  \[
    A \text{ Artinian}
    \Longleftrightarrow
    A \text{ Noetherian 且 } \dim A=0.
  \]
\end{itemize}

\paragraph{结构定理}
Artinian 环唯一分解成有限个 Artinian 局部环的直积。

\paragraph{局部情形}
Noether 局部环满足
\[
\mathfrak m^n=0 \text{ for some } n
\Longleftrightarrow
A \text{ Artinian}.
\]

\paragraph{考前提示}
\begin{itemize}
  \item “Artinian = 零维 Noetherian” 是高频口号。
  \item Artinian 局部环中，每个元素不是单位就是幂零元。
\end{itemize}

\section{第6章：Associated primes and Primary decomposition}

\subsection{6.1 Associated primes}

\paragraph{定义}
\[
\operatorname{Ass}(M)=\{\operatorname{ann}(x)\mid x\in M\}.
\]

\paragraph{关系}
\begin{itemize}
  \item \(\operatorname{Ass}(M)\subset \operatorname{Supp}(M)\)。
  \item \(\operatorname{Supp}(M)\) 的极小元就是 associated primes 中的极小元。
  \item 对有限生成模，\(\operatorname{Ass}(M)\) 是有限集。
\end{itemize}

\paragraph{为什么重要}
associated primes 记录了模的“零因子来源”和不可约几何成分。

\subsection{6.2 Primary decomposition}

\paragraph{定义}
子模或理想可写成若干 primary 子模（或 primary ideals）的交。

\paragraph{最重要的结论}
\begin{itemize}
  \item Noether 环上的有限生成模都有 primary decomposition。
  \item 不可约冗余去掉后，出现的根理想集合是唯一的。
  \item 极小 associated primes 对应的 primary 成分是唯一决定的。
\end{itemize}

\paragraph{考前提示}
\begin{itemize}
  \item 记住“存在性靠 Noether，唯一性只保留到某种有限程度”。
  \item primary decomposition 是后面研究高度、unmixed、CM 性质的基础。
\end{itemize}

\section{第7章：One-dimensional integrally closed Noetherian domains}

\subsection{7.1 DVR and Dedekind domains}

\paragraph{DVR}
离散赋值环是最基本的一维局部模型。对 Noether 局部整环 \(A\)，一维且整闭时会导向 DVR。

\paragraph{Dedekind domain}
Dedekind 整环定义为：
\begin{itemize}
  \item Noetherian；
  \item \(\dim A=1\)；
  \item integrally closed。
\end{itemize}

\paragraph{核心结论}
\begin{itemize}
  \item Dedekind 整环的每个非零素理想都是极大理想。
  \item 每个非零理想唯一分解为素理想幂的乘积。
  \item 非零素理想局部化后得到 DVR。
\end{itemize}

\paragraph{例子}
\begin{itemize}
  \item PID 是 Dedekind 整环。
  \item 数域整数环是 Dedekind 整环。
\end{itemize}

\subsection{7.2 Fractional ideals}

\paragraph{关键观点}
Dedekind 整环可以用分式理想完全刻画。

\paragraph{必须记住的判别}
\begin{itemize}
  \item 局部整环 \(A\) 是 DVR
  \(\Longleftrightarrow\)
  每个非零分式理想可逆。
  \item 整环 \(A\) 是 Dedekind
  \(\Longleftrightarrow\)
  每个非零分式理想可逆。
\end{itemize}

\paragraph{类群}
Dedekind 整环的非零分式理想模去主分式理想，得到 ideal class group。
要记住：
\[
\operatorname{Cl}(A)=0
\Longleftrightarrow
A \text{ 是 PID}
\Longleftrightarrow
A \text{ 是 UFD}.
\]

\paragraph{考前提示}
第七章的主旨是：一维整闭 Noetherian 整环的结构已经非常刚性，几乎完全由理想理论控制。

\section{第8章：Graded rings and Completions}

\subsection{8.1 Graded rings and modules}

\paragraph{核心作用}
分次结构是研究 Hilbert 函数、associated graded ring、Rees 代数的入口。

\paragraph{要点}
\begin{itemize}
  \item 分次 Noether 性可以转化成 \(A_0\) 的 Noether 性与有限生成性。
  \item 后面会把局部环上的 filtration 转成分次对象来研究。
\end{itemize}

\subsection{8.2 Artin--Rees lemma}

\paragraph{内容}
若 \(A\) Noetherian，\(I\subset A\) 理想，\(N\subset M\) 有限生成模，则 \(I\)-adic filtration 与子模交的行为最终稳定。

\paragraph{意义}
\begin{itemize}
  \item 它控制 \(N\cap I^nM\) 的形状。
  \item 它是完备化正合性、局部 flat 判别、Hilbert 多项式理论的关键技术工具。
\end{itemize}

\subsection{8.3 Local criterion of flatness}

\paragraph{主旨}
在 Noether 局部环境中，flatness 可用模掉 \(\mathfrak m\) 后的信息和 \(\operatorname{Tor}\) 消失来判别。

\paragraph{考前提示}
这节常和 Nakayama、Artin--Rees 一起联动，是同调控制局部性质的典型例子。

\subsection{8.4 Topology and completion}

\paragraph{\(I\)-adic topology}
理想幂
\[
I^0\supset I\supset I^2\supset \cdots
\]
定义 \(I\)-adic topology，其完备化是
\[
\widehat A=\varprojlim A/I^n.
\]

\paragraph{要点}
\begin{itemize}
  \item 完备化是函子性的。
  \item 在 Noether 情形，有限生成模的完备化与正合列兼容。
  \item 完备化后的局部环仍是局部环。
\end{itemize}

\subsection{8.5 Applications to Noetherian rings}

\paragraph{高频结论}
\begin{itemize}
  \item Noether 环的完备化仍是 Noether 环。
  \item associated graded ring 和 completion 紧密相关。
  \item 稳定 filtration 与 \(I\)-adic topology 给出同样的渐近信息。
\end{itemize}

\paragraph{考前提示}
第八章是第九章维数理论的技术准备章。

\section{第9章：Dimension theory of Noetherian rings}

\subsection{总框架}

\paragraph{定义}
\begin{itemize}
  \item
  \[
    \dim A=\sup\{\text{素理想严格链长度}\}.
  \]
  \item 对 \(\mathfrak p\in\operatorname{Spec}A\),
  \[
    \operatorname{ht}\mathfrak p=\dim A_{\mathfrak p}.
  \]
  \item 对理想 \(I\),
  \[
    \operatorname{ht}I=\inf_{\mathfrak p\supset I}\operatorname{ht}\mathfrak p.
  \]
\end{itemize}

\paragraph{基本不等式}
\[
\operatorname{ht}I+\dim A/I\le \dim A.
\]

\subsection{9.1 Hilbert functions}

\paragraph{主线}
对分次环和分次模，Hilbert 级数是有理函数；当生成元次数为 \(1\) 时，长度函数最终由多项式给出。

\paragraph{最重要信息}
\begin{itemize}
  \item Hilbert 多项式的次数记录了维数信息。
  \item 非零因子模掉一次，会把相关次数降 \(1\)。
\end{itemize}

\subsection{9.2 Noether 局部环的维数理论}

\paragraph{核心定理}
设 \((A,\mathfrak m)\) 为 Noether 局部环，则
\[
\dim A=d(A)=\nu(A),
\]
其中 \(d(A)\) 来自 Hilbert 多项式次数，\(\nu(A)\) 是 \(\mathfrak m\)-primary 理想最少生成元个数（III.Thm 9.2.7）。

\paragraph{系统参数}
若 \(\dim A=d\)，生成 \(\mathfrak m\)-primary 理想的一组 \(d\) 个元素称为 system of parameters。

\paragraph{高频推论}
\begin{itemize}
  \item
  \[
    \dim A\le \dim_k(\mathfrak m/\mathfrak m^2).
  \]
  \item 若 \(I=(x_1,\dots,x_r)\)，则
  \[
    \operatorname{ht}I\le r.
  \]
  \item 若 \(x\in \mathfrak m\) 是非零因子，则
  \[
    \dim A/(x)=\dim A-1.
  \]
  \item 完备化不改变维数。
\end{itemize}

\subsection{9.3 Transcendental dimension}

\paragraph{维数公式}
若 \(A\) 是有限生成 \(k\)-代数整环，\(K=\operatorname{Frac}(A)\)，则
\[
\dim A=\operatorname{tr.deg}_k K.
\]

\paragraph{意义}
这是代数几何维数与交换代数维数统一的关键点。

\paragraph{考前速记}
\begin{itemize}
  \item 局部环维数由参数组长度控制。
  \item Hilbert 多项式次数就是维数。
  \item 高度估计 \(\operatorname{ht}(x_1,\dots,x_r)\le r\) 必须熟。
\end{itemize}

\section{第10章：Cohen--Macaulay, Normal, Regular}

\subsection{10.1 Regular local rings}

\paragraph{定义}
对 Noether 局部环 \((A,\mathfrak m,k)\)，若
\[
\dim A=\dim_k(\mathfrak m/\mathfrak m^2),
\]
则 \(A\) 为 regular local ring。

\paragraph{关键判别}
\[
A \text{ regular}
\Longleftrightarrow
G_{\mathfrak m}(A)\cong k[t_1,\dots,t_d].
\]

\paragraph{直接后果}
\begin{itemize}
  \item 完备化保持 regular。
  \item \(k[[X_1,\dots,X_d]]\) 是 regular local ring。
  \item DVR 是一维 regular local ring。
  \item regular local ring 是整环，并且整闭，事实上还是 UFD。
\end{itemize}

\subsection{10.2 Regular sequences, Koszul complex, depth}

\paragraph{正则列}
\(a_1,\dots,a_r\) 是 \(M\)-regular sequence，表示每一步都不是前一步商模上的零因子，并且最终没有把模整个杀掉。

\paragraph{depth}
\[
\operatorname{depth}_I M
=
\text{理想 } I \text{ 中最长 } M\text{-regular sequence 的长度}.
\]

\paragraph{Koszul 复形}
Koszul 复形是检测正则列的显式工具。若 \(x_1,\dots,x_n\) 是 \(M\)-regular，则
\[
H_p(x,M)=0\quad (p>0),\qquad H_0(x,M)=M/xM.
\]

\paragraph{考前提示}
\begin{itemize}
  \item 正则列是“没有隐藏零因子”的序列。
  \item depth 衡量的是模的深度，而不是维数。
  \item 以后 CM 的定义就是 depth 和维数相等。
\end{itemize}

\subsection{10.3 Cohen--Macaulay rings}

\paragraph{定义}
对有限生成非零模 \(M\)，若
\[
\operatorname{depth}M=\dim M,
\]
则称 \(M\) 是 Cohen--Macaulay。

\paragraph{关键结论}
\begin{itemize}
  \item regular sequence 模掉后，CM 性保持。
  \item 局部化保持 CM。
  \item regular local ring 一定是 CM。
  \item 在 CM 局部环中，高度、depth、参数组、正则列之间高度统一。
\end{itemize}

\paragraph{必须会的口号}
\[
\text{Cohen--Macaulay}=\text{depth 尽可能大}.
\]

\subsection{10.4 Normal rings}

\paragraph{定义}
Noether 环 \(A\) 是 normal，如果每个 \(A_{\mathfrak p}\) 都是整闭整环。

\paragraph{一维情形}
一维 Noether 局部环中，
\[
\text{regular}
\Longleftrightarrow
\text{normal}
\Longleftrightarrow
\text{DVR}.
\]

\paragraph{Serre 判别}
对 Noether 环，
\[
A \text{ normal}
\Longleftrightarrow
(S_2)+(R_1).
\]

\paragraph{考前理解}
\begin{itemize}
  \item \((R_1)\) 管余维 \(1\) 处的 regular 性。
  \item \((S_2)\) 管整体深度不能太差。
\end{itemize}

\subsection{10.5 Homological dimension and Regular rings}

\paragraph{全局维数}
\[
\operatorname{gl.dim}A=\sup_M \operatorname{proj.dim}M.
\]

\paragraph{局部判别}
对 Noether 局部环 \((A,\mathfrak m,k)\)，有
\[
\operatorname{gl.dim}A=\operatorname{proj.dim}_A k.
\]

\paragraph{Serre 对 regular 的同调刻画}
\[
A \text{ regular}
\Longleftrightarrow
\operatorname{gl.dim}A=\dim A
\Longleftrightarrow
\operatorname{gl.dim}A<\infty.
\]

\paragraph{Auslander--Buchsbaum}
若 \(M\) 有有限投射维数，则
\[
\operatorname{proj.dim}M+\operatorname{depth}M=\operatorname{depth}A.
\]

\paragraph{考前提示}
regular 不只是“极大理想生成元个数刚好等于维数”，也等价于“同调维数有限”。

\section*{全书关系图}
\phantomsection
\addcontentsline{toc}{section}{全书关系图}

\[
\text{localization} \Rightarrow \text{把问题化到局部环};
\]
\[
\text{flat} \Rightarrow \text{控制张量与 going-down};
\]
\[
\text{integral} \Rightarrow \text{控制谱映射与 going-up};
\]
\[
\text{Noetherian} \Rightarrow \text{有限性、primary decomposition、Artin--Rees};
\]
\[
\text{completion + graded} \Rightarrow \text{Hilbert 函数与维数理论};
\]
\[
\text{dimension} \Rightarrow \text{参数组、高度估计、超越次数};
\]
\[
\text{regular} \Rightarrow \text{Cohen--Macaulay and normal}.
\]

更具体地：
\[
\text{regular} \Rightarrow \text{Cohen--Macaulay} \Rightarrow (S_k)\ \text{行为良好},
\]
\[
\text{regular} \Rightarrow \text{normal},
\qquad
\text{normal} \Longleftrightarrow (S_2)+(R_1)\ \text{(Noether)}.
\]

\section*{最后的考前清单}
\phantomsection
\addcontentsline{toc}{section}{最后的考前清单}

\begin{itemize}
  \item 会写出 \(\operatorname{Spec}A\)、\(V(I)\)、\(\sqrt I\)、\(\operatorname{nilrad}(A)\)。
  \item 会用 Nakayama 引理判断最少生成元个数。
  \item 会写出局部化后谱和模的变化。
  \item 知道 flat 的三个核心标签：张量正合、\(\operatorname{Tor}_1=0\)、可局部检验。
  \item 知道 integral 的三个核心标签：首一方程、谱满射、going-up。
  \item 会说 Hilbert basis theorem。
  \item 知道 Artinian \(=\) 零维 Noetherian。
  \item 知道 associated primes 和 primary decomposition 的作用。
  \item 知道 Dedekind 整环等价于“一维 Noetherian 整闭”，且非零理想唯一分解。
  \item 知道完备化与 Artin--Rees 是第九章的技术准备。
  \item 会背 Krull 维数定理的局部版本：\(\dim A=d(A)=\nu(A)\)。
  \item 会用高度估计 \(\operatorname{ht}(x_1,\dots,x_r)\le r\)。
  \item 知道 regular sequence、depth、Koszul complex 三者关系。
  \item 知道 CM 的定义是 depth \(=\) dim。
  \item 知道 normal 的 Serre 判别和 regular 的 Serre 同调判别。
\end{itemize}

\end{document}
